\documentclass[10pt,twocolumn]{article}

% ── Core packages ──────────────────────────────────────────────────────────
\usepackage[T1]{fontenc}
\usepackage[utf8]{inputenc}
\usepackage[margin=0.75in, top=0.85in, bottom=0.85in]{geometry}
\usepackage{times}
\usepackage{graphicx}
\usepackage{booktabs}
\usepackage{array}
\usepackage{tabularx}
\usepackage{microtype}
\usepackage{enumitem}
\usepackage{caption}
\usepackage{fancyhdr}
\usepackage{titlesec}
\usepackage{parskip}
\usepackage{hyperref}
\usepackage{amsmath}
\usepackage{multicol}
\usepackage{balance}

% ── IEEE-like formatting ──────────────────────────────────────────────────
\setlength{\columnsep}{0.25in}
\setlength{\parindent}{1em}
\setlength{\parskip}{0pt}

% Section: centered, bold, all caps, Roman numeral
\titleformat{\section}[block]
  {\normalfont\bfseries\centering\normalsize}{}{0em}
  {\thesection.\enspace\MakeUppercase}
\titlespacing{\section}{0pt}{8pt plus 2pt}{4pt plus 1pt}

% Subsection: left-aligned, italic, letter
\titleformat{\subsection}[block]
  {\normalfont\itshape\normalsize}{}{0em}
  {\Alph{subsection}.\enspace}
\titlespacing{\subsection}{0pt}{6pt plus 1pt}{2pt}

% Sub-subsection: run-in italic bold
\titleformat{\subsubsection}[runin]
  {\normalfont\itshape\normalsize\bfseries}{}{0em}{}[.]
\titlespacing{\subsubsection}{0pt}{4pt}{0.5em}

% Abstract box
\newenvironment{ieeeabstract}{%
  \begin{center}
  \begin{minipage}{0.9\linewidth}
  \small\noindent\textbf{Abstract---}\ignorespaces}
 {\end{minipage}\end{center}\vspace{4pt}}

% Header / footer
\pagestyle{fancy}
\fancyhf{}
\fancyhead[C]{\footnotesize Systems Analysis \& Design --- Workshop No.~2, Semester 2026-I}
\fancyfoot[C]{\footnotesize\thepage}
\renewcommand{\headrulewidth}{0.4pt}

% Caption style
\captionsetup{font=footnotesize,labelfont=bf,labelsep=period}

% Column type for long text in tables
\newcolumntype{L}[1]{>{\raggedright\arraybackslash\footnotesize}p{#1}}

% References style
\renewcommand{\refname}{\normalfont\bfseries\centering\normalsize REFERENCES}

\begin{document}

% ── Title block (single column) ────────────────────────────────────────────
\twocolumn[{%
\begin{center}
  {\large\bfseries Community Security Alert System (CSAS):\\[4pt]
   Systems Design --- Workshop No.~2}\\[10pt]
  {\normalsize
   Felipe~Jose~Garzon~Herrera,~\textit{20251020132},\quad
   Juan~Esteban~Quintero~Gordillo,~\textit{20251020137},\\[2pt]
   Henry~Samuel~Garrido~Medina,~\textit{20251020125},\quad
   Gabriel~Mateo~Cusba~Marin,~\textit{20251020128}\\[6pt]
   \textit{Systems Analysis \& Design, Semester 2026-I}\\
   \textit{Eng.\ Carlos Andr\'es Sierra, M.Sc.}\\
   \textit{Computer Engineering Program, Universidad Distrital Francisco
   Jos\'e de Caldas, Bogot\'a D.C., Colombia}}\\[10pt]
\end{center}
\begin{ieeeabstract}
This document presents the systems design specification for the Community
Security Alert System (CSAS), a crowd-sourced mobile and web platform
engineered to improve safety at Universidad Distrital Francisco Jos\'e de
Caldas. The design builds directly on the empirical findings of Workshop
No.~1, which identified three critical gaps in the current campus security
posture: a persistent under-reporting culture, a spatial and temporal
concentration of incidents in specific zones and hours, and the absence
of a centralized real-time reporting tool. The proposed architecture
adopts a microservices-based, event-driven approach organized across five
layers. Key design decisions include a hybrid AI-and-human verification
pipeline, geofenced alert dispatch with zone-weighted and time-weighted
priority logic, and a community engagement module targeting the adoption
gap. By integrating a mobile-first reporting interface with automated
verification and a multi-channel notification engine, the system is
designed to reduce average emergency response times from 11.4 minutes to
under 5 minutes.

\textbf{Index Terms---}campus security, crowd-sourced incident
reporting, microservices architecture, real-time notifications,
geofenced alerts, verification pipeline, event-driven design.
\end{ieeeabstract}
\vspace{8pt}
}]

% =============================================================================
\section{Introduction}
% =============================================================================

The CSAS project addresses a systemic failure in traditional security
reporting at Universidad Distrital Francisco Jos\'e de Caldas. Workshop
No.~1 established, through six independent data collection methods, that
the current campus security posture suffers from three compounding
problems. First, most incidents go unreported through official channels
because the friction of current reporting tools is too high. Second,
incidents are spatially and temporally concentrated: all observed
suspicious events occurred during evening or night hours, and the parking
lot was consistently identified as the highest-risk zone. Third, there
is no centralized, real-time incident management tool; security staff
rely on phone calls and walk-up reports, which introduces communication
delays averaging~11.4 minutes per incident.

The goal of this design phase is to translate those analytical findings
into a coherent technical architecture that directly addresses each
identified gap. Table~\ref{tab:gaps} summarizes the operational
difference between the current state and the design targets established
for this workshop.

\begin{table}[ht]
\caption{Operational Gap Analysis: Current State vs.\ Design Targets}
\label{tab:gaps}
\renewcommand{\arraystretch}{1.15}
\centering
\begin{tabular}{L{1.8cm} L{1.5cm} L{1.4cm} L{1.1cm}}
\toprule
\textbf{Dimension} & \textbf{Baseline} & \textbf{Target} & \textbf{Gap} \\
\midrule
Response time      & 11.4 min   & $<$5 min    & $-$6.4\,min \\
Reporting rate     & 41\%       & $>$85\%     & $+$44\,pp   \\
User adoption      & $\sim$12\% & $>$60\%     & $+$48\,pp   \\
Delivery success   & $\sim$87\% & $>$99.9\%   & $+$12.9\,pp \\
\bottomrule
\end{tabular}
\end{table}

The remainder of this document is organized as follows. Section~II
defines functional and non-functional requirements. Section~III
describes the system architecture. Section~IV presents the
implementation strategy. Section~V addresses risk management.
Section~VI defines performance metrics and testing. Section~VII
discusses ethical considerations. Section~VIII concludes.

% =============================================================================
\section{Requirements Analysis}
% =============================================================================

\subsection{Functional Requirements}

The following requirements are derived directly from Workshop~No.~1.
Each is traceable to at least one empirical finding.

\begin{itemize}[leftmargin=*, itemsep=1pt, topsep=2pt]
\item \textbf{FR-01 (Incident Reporting):} Users shall report incidents
by specifying type, GPS location, and description. Addresses the
under-reporting barrier.
\item \textbf{FR-02 (Geofenced Notifications):} Real-time alerts shall
be broadcast to users within a 500\,m radius of a verified incident.
Responds to the spatial concentration of incidents in high-risk zones.
\item \textbf{FR-03 (AI Verification):} An AI plausibility scorer shall
perform initial metadata validation and noise filtering. Mitigates data
fragmentation in current manual records.
\item \textbf{FR-04 (Administrative Dashboard):} Security personnel
shall have a web interface to review, escalate, or resolve active alerts.
Replaces fragmented phone and walk-up reporting.
\item \textbf{FR-05 (Crowd-Sourced Validation):} Users shall confirm or
flag reports submitted by others. Leverages the high
willingness-to-participate rate observed in the student survey.
\item \textbf{FR-06 (Automatic Geolocation):} The mobile client shall
automatically capture GPS coordinates for every submitted report. Based
on spatial precision needs identified in direct observation sessions.
\item \textbf{FR-07 (Centralized Database):} A persistent incident log
shall be maintained for historical pattern analysis. Addresses the
current state in which incidents go unrecorded.
\item \textbf{FR-08 (Emergency SOS):} A one-touch emergency button shall
trigger immediate security dispatch. Targets a reduction in response
time of more than half the current baseline.
\item \textbf{FR-09 (Analytics Engine):} The system shall generate
periodic security heatmaps for university management. Provides
data-driven oversight currently absent.
\item \textbf{FR-10 (Media Attachments):} Users shall attach image or
video evidence to reports to improve verification quality.
\end{itemize}

\subsection{Non-Functional Requirements}

\begin{itemize}[leftmargin=*, itemsep=1pt, topsep=2pt]
\item \textbf{NFR-01 (Latency):} End-to-end latency from report
submission to alert broadcast must be under~30\,s for~95\% of incidents.
\item \textbf{NFR-02 (Reliability):} The notification engine shall
maintain a~99.9\% delivery success rate across all channels.
\item \textbf{NFR-03 (Availability):} The system must maintain~99.9\%
uptime, with guaranteed operation during high-risk hours (18:00--22:00).
\item \textbf{NFR-04 (Privacy/Security):} All data must be encrypted
via TLS~1.3 and comply with Ley~1581 de~2012.
\item \textbf{NFR-05 (Scalability):} The API gateway must handle
a~300\% surge in concurrent requests without performance loss.
\item \textbf{NFR-06 (Usability):} The reporting workflow must require
no more than three interactions to submit an alert.
\item \textbf{NFR-07 (Maintainability):} Any microservice shall be
updatable or replaceable independently.
\item \textbf{NFR-08 (Data Integrity):} No alert shall broadcast as
verified without at least two community confirmations or one
administrative approval.
\end{itemize}

% =============================================================================
\section{System Architecture}
% =============================================================================

\subsection{Architectural Style}

The CSAS adopts a microservices-based architecture in which the system
is decomposed into independent components that communicate through
well-defined interfaces. This choice is driven by three characteristics
of the problem: the verification, notification, and analytics subsystems
have independent load profiles and scaling needs; alert dispatch must be
decoupled from incident submission to prevent cascading delays during
high-volume periods; and the system must integrate with external actors
through stable, versioned interfaces.

\subsection{System Layers}

The architecture is organized across five layers:

\begin{itemize}[leftmargin=*, itemsep=1pt, topsep=2pt]
\item \textbf{Presentation:} Mobile application and web portal for
incident submission and alert reception.
\item \textbf{Application:} API gateway handling authentication, rate
limiting, and microservice routing.
\item \textbf{Service:} Five domain microservices (Section~III-C).
\item \textbf{Data:} Persistent storage for incidents, user profiles,
media, audit logs, and analytics, backed by a Redis cache for
low-latency reads.
\item \textbf{Integration:} Interfaces to SIURE~UD, local authorities,
and the SMS gateway.
\end{itemize}

\subsection{Microservices Design}

Table~\ref{tab:services} describes the five core microservices. Each
publishes and consumes domain events over the message broker (RabbitMQ),
ensuring loose coupling across the service layer.

\begin{table}[ht]
\caption{Core Microservices and Responsibilities}
\label{tab:services}
\renewcommand{\arraystretch}{1.15}
\centering
\begin{tabular}{L{1.5cm} L{5.5cm}}
\toprule
\textbf{Service} & \textbf{Responsibility} \\
\midrule
Incident & Receives, validates, and persists reports. Enriches each
event with geospatial metadata and publishes \textit{IncidentCreated}. \\
\addlinespace
Verification & Orchestrates the verification workflow using an adaptive
threshold algorithm and AI plausibility scorer. Supports multi-user
confirmation, administrative override, and priority bypass for
high-severity types. Publishes \textit{IncidentVerified}. \\
\addlinespace
Dispatcher & Computes the notification target set using zone-weighted
and time-weighted priority logic. Distributes alerts via push and SMS
fallback. \\
\addlinespace
User & Manages registration, roles, and notification preferences.
Implements alert fatigue mitigation by dampening non-critical alerts for
high-dismissal users. \\
\addlinespace
Analytics & Consumes domain events to populate the data warehouse.
Produces security dashboards and periodic heatmap reports (FR-09). \\
\bottomrule
\end{tabular}
\end{table}

\subsection{External Integrations}

\begin{table}[ht]
\caption{External Integrations}
\label{tab:integrations}
\renewcommand{\arraystretch}{1.15}
\centering
\begin{tabular}{L{1.4cm} L{1.1cm} L{4.5cm}}
\toprule
\textbf{System} & \textbf{Type} & \textbf{Description} \\
\midrule
SIURE~UD & REST API & Bidirectional bridge: CSAS forwards verified
incidents; SIURE~UD provides historical baseline. Complements the
existing platform without replacing it. \\
\addlinespace
Authorities & REST API & Secure endpoint activated for high-severity
incidents requiring external intervention. Modeled after the Polic\'ia
Nacional / Vigila~App integration. \\
\addlinespace
SMS Gateway & Ext.\ service & Fallback channel for users without push
capability (NFR-02). Auto-activates after a 10-second push delivery
timeout. \\
\bottomrule
\end{tabular}
\end{table}

Figs.~\ref{fig:architecture} and~\ref{fig:processflow} provide visual
representations of the architecture and main process flow described in
Sections~III-B through~III-D.

% ── Figures (place fig1_architecture.png and fig2_processflow.png
%    in the same folder as this .tex file before compiling) ─────────────────
\begin{figure}[!ht]
\centering
\includegraphics[width=\columnwidth]{fig1_architecture.png}
\caption{CSAS System Architecture Diagram. The REST API gateway acts as
the central connectivity node between the presentation, service, and
data layers.}
\label{fig:architecture}
\end{figure}

\begin{figure}[!ht]
\centering
\includegraphics[width=\columnwidth]{fig2_processflow.png}
\caption{Main Process Flow Diagram. Information flows from user
submission through the verification pipeline to alert dispatch and
database logging.}
\label{fig:processflow}
\end{figure}

% =============================================================================
\section{Technical Implementation Strategy}
% =============================================================================

\subsection{Technology Stack}

Table~\ref{tab:stack} presents the recommended technology stack. Each
choice is justified with reference to the requirements defined in
Section~II.

\begin{table}[ht]
\caption{Recommended Technology Stack}
\label{tab:stack}
\renewcommand{\arraystretch}{1.15}
\centering
\begin{tabular}{L{1.6cm} L{1.8cm} L{3.6cm}}
\toprule
\textbf{Component} & \textbf{Technology} & \textbf{Justification} \\
\midrule
Mobile app & React Native & Single codebase for Android/iOS; supports
$\leq$3-interaction reporting (NFR-06). \\
\addlinespace
API gateway & Node.js / Nginx & High I/O throughput for concurrent
submissions; handles routing and rate limiting. \\
\addlinespace
Backend services & Node.js & Asynchronous, event-driven I/O supports
NFR-01 ($<$30\,s latency). \\
\addlinespace
Verification Engine & Python / FastAPI & Python's ML ecosystem
simplifies integration of the AI plausibility scorer. \\
\addlinespace
Message broker & RabbitMQ & Priority queue support; decouples ingestion
from processing to mitigate volume spikes. \\
\addlinespace
Database & PostgreSQL + PostGIS & Native geospatial queries for
zone-weighted routing and heatmaps (FR-06, FR-09). \\
\addlinespace
Push notifications & Firebase FCM & Cross-platform delivery; supports
99.9\% delivery target (NFR-02). \\
\addlinespace
SMS fallback & Twilio API & Guaranteed delivery for users without
smartphone access; supports equity of access. \\
\addlinespace
Infrastructure & AWS / Azure & Container orchestration via Kubernetes
enables auto-scaling (NFR-03, NFR-05). \\
\bottomrule
\end{tabular}
\end{table}

\subsection{Implementation Phases}

\subsubsection{Phase 1 --- Planning and Architecture (Month~1)}
Definition of the final architecture, validation of all requirements,
and development environment setup. \textit{Success criterion:} all FR
and NFR are traceable to Workshop~No.~1 findings. \textit{Deliverable:}
architecture diagrams and requirements documentation.

\subsubsection{Phase 2 --- Development (Month~2)}
Implementation of frontend, backend microservices, database schema, and
messaging infrastructure. \textit{Success criterion:} end-to-end
prototype in which a user can submit a report (FR-01) and receive a
geofenced alert (FR-02). \textit{Deliverable:} functional prototype.

\subsubsection{Phase 3 --- Testing and Validation (Month~3)}
Execution of the full testing strategy (Section~VI-B). \textit{Success
criterion:} system meets NFR-01 ($<$30\,s) and NFR-02 (99.9\%) under
simulated load. \textit{Deliverable:} test reports.

\subsubsection{Phase 4 --- Deployment and Monitoring (Month~4)}
Cloud deployment, external integrations, and production monitoring
activation. \textit{Success criterion:} 99.9\% uptime under real-world
conditions. \textit{Deliverable:} deployed system with live dashboard.

% =============================================================================
\section{Risk and Complexity Management}
% =============================================================================

\subsection{Sensitivity Parameter Mitigation}

Workshop~No.~1 identified five operational parameters with the strongest
influence on CSAS behavior. Left unaddressed, each represents a distinct
failure path. Table~\ref{tab:sensitivity} maps each parameter to its
design mechanism.

\begin{table*}[ht]
\caption{Sensitivity Parameter Risk and Design Mitigation}
\label{tab:sensitivity}
\renewcommand{\arraystretch}{1.2}
\centering
\begin{tabular}{L{2.4cm} L{4.0cm} L{9.5cm}}
\toprule
\textbf{Parameter} & \textbf{Risk} & \textbf{Design Mechanism} \\
\midrule
Report volume spike
& Verification engine saturates; alert latency exceeds targets.
& RabbitMQ decouples report ingestion from verification processing.
  Auto-scaling policies on the Verification Engine pods. Priority
  queue lanes for high-severity events that bypass the standard queue. \\
\addlinespace
Low user participation
& Real incidents go undetected; system loses credibility.
& Low-friction reporting interface (FR-01, $\leq$3 steps). Anonymous
  reporting option. Community engagement module (Phase~3). Integration
  with SIURE~UD as proof of institutional endorsement. \\
\addlinespace
Night-time concentration
& Higher alert volume during off-hours; fewer moderators available.
& Time-weighted priority logic elevates night-hour alert priority.
  On-call escalation path notifies the on-duty security team when no
  active moderators are available. \\
\addlinespace
High-risk zone clustering
& Alert storms in parking lot and dormitory zones saturate users.
& Zone-based rate limiting on alert broadcast. Spatial clustering
  algorithm aggregates co-located simultaneous reports into a single
  composite alert to prevent redundancy. \\
\addlinespace
Adverse weather conditions
& Campus movement decreases; visibility drops; incident risk rises.
& Weather API integration feeds real-time conditions into the risk
  scoring model. Adverse weather triggers automatic elevation of
  monitoring sensitivity in high-risk zones. \\
\bottomrule
\end{tabular}
\end{table*}

\subsection{Feedback Loops and Emergent Behavior}

\subsubsection{Reinforcing loop}
The system exhibits a positive reinforcing dynamic: useful alerts
increase user vigilance, which increases reporting, which generates more
alerts, progressively improving system effectiveness as adoption
grows~\cite{meadows2008}. The community engagement module (Phase~3)
targets the critical adoption threshold of approximately 20--25\% of the
campus population, the point at which report density sustains the
verification mechanism without heavy reliance on security staff.

\subsubsection{Balancing loop}
A counterbalancing dynamic operates in the opposite direction: excessive
alerts cause users to ignore them (alert fatigue), reducing system
effectiveness over time~\cite{meadows2008}. The alert fatigue mitigation
module in the User Service dampens non-critical notifications for users
with high dismissal rates. Simultaneously, the Verification Engine's
quality filters prevent low-confidence reports from reaching the
dispatcher. Calibrating the right confirmation threshold --- fast enough
for genuine emergencies, strict enough to suppress noise --- is the most
critical design tension in the system.

Two non-linear effects are also worth noting: the relationship between
adoption and report volume is non-linear (a small user increase can
produce a disproportionate volume surge), and alert activity is not
uniformly distributed across campus zones or time of day. For this
reason, the dispatcher applies zone-weighted and time-weighted priority
logic rather than treating all reports equally~\cite{checkland1981}.

\subsection{Failure Mode Analysis}

Table~\ref{tab:failures} maps each major architectural dependency to a
failure mode, its impact, and the corresponding mitigation.

\begin{table*}[ht]
\caption{Failure Mode Analysis and Mitigation Strategies}
\label{tab:failures}
\renewcommand{\arraystretch}{1.2}
\centering
\begin{tabular}{L{2.8cm} L{4.0cm} L{9.2cm}}
\toprule
\textbf{Failure Mode} & \textbf{Impact} & \textbf{Mitigation} \\
\midrule
API Gateway outage
& All user submissions fail; system is unreachable.
& Multi-region deployment with health checks and automatic failover.
  Mobile app queues reports locally and retries on reconnection. \\
\addlinespace
Verification Engine overload
& Alert latency exceeds NFR-01.
& Auto-scaling on Verification Engine pods. High-severity incidents
  bypass AI scoring via a priority dispatch path. \\
\addlinespace
Push notification failure
& Users do not receive critical alerts.
& SMS fallback via Twilio activates after a 10-second delivery
  timeout. Delivery receipts confirm end-to-end success. \\
\addlinespace
Database unavailability
& Risk of incident data loss.
& PostgreSQL primary-replica replication. Automated daily backups.
  Write-ahead logging ensures no committed data is lost. \\
\addlinespace
AI scorer false positive
& Legitimate incidents blocked before dispatch.
& Human moderator review queue for all low-confidence scores.
  Configurable override timer (default: 90\,s) escalates to human
  review. \\
\addlinespace
Mobile connectivity dead zone
& Submissions fail in low-coverage campus areas.
& Offline queuing in the mobile app. SMS submission as an
  alternative channel. A connectivity audit is recommended prior to
  Phase~1 infrastructure planning. \\
\bottomrule
\end{tabular}
\end{table*}

% =============================================================================
\section{Performance and Quality Assurance}
% =============================================================================

\subsection{Key Performance Indicators}

Table~\ref{tab:kpis} defines the measurable targets against which the
implemented system will be evaluated. Each KPI is aligned with a
specific requirement from Section~II.

\begin{table}[ht]
\caption{Key Performance Indicators}
\label{tab:kpis}
\renewcommand{\arraystretch}{1.15}
\centering
\begin{tabular}{L{2.2cm} L{1.8cm} L{1.5cm} L{1.0cm}}
\toprule
\textbf{KPI} & \textbf{Target} & \textbf{Method} & \textbf{Freq.} \\
\midrule
Alert latency (end-to-end) & $<$30\,s (NFR-01)       & System logs & Daily   \\
Delivery success rate       & $>$99.9\% (NFR-02)      & Notif.\ logs & Daily  \\
System availability         & 99.9\% (NFR-03)         & Monitoring  & Monthly \\
Incident reporting rate     & $>$85\% logged          & DB records  & Weekly  \\
User adoption rate          & $>$60\% of students     & App analytics & Monthly \\
Verification latency        & $<$60\,s                & Event logs  & Daily   \\
Alert accuracy              & $\geq$2 confirm.\ or 1 admin & Validation & Daily \\
Scalability under surge     & 300\% surge (NFR-05)    & Load tests  & Monthly \\
\bottomrule
\end{tabular}
\end{table}

\subsection{Testing Strategy}

\subsubsection{Unit testing}
Validates individual components in each microservice. Each service shall
maintain at least~80\% code coverage.

\subsubsection{Integration testing}
Verifies the full event chain: Incident Service $\rightarrow$
Verification Engine $\rightarrow$ Alert Dispatcher. Tests run against
a staging environment with seeded campus zone data.

\subsubsection{Load testing}
Evaluates system behavior under a 300\% traffic surge (NFR-05).
Target: zero timeout errors and latency below~30\,s at the~95th
percentile under full load.

\subsubsection{User acceptance testing (UAT)}
A pilot group of students and security staff validates usability and
correctness under realistic conditions. Success criterion: at least 85\%
of participants complete a report submission in under two minutes.

\subsubsection{Security testing}
Verifies compliance with NFR-04 and Ley~1581 de~2012: data encryption
(TLS~1.3), secure authentication, and access control mechanisms.

% =============================================================================
\section{Ethical and Safety Considerations}
% =============================================================================

\subsection{Privacy and Data Protection}

The CSAS collects three categories of personal data: registration
information (name, institutional email, and campus role), real-time
location data generated during incident submission, and incident media
including photographs and video. Each category carries distinct privacy
obligations under Colombian law.

All data handling must comply with Ley~1581 de~2012, Colombia's general
personal data protection regime~\cite{ley1581}. Four obligations are
directly binding on the CSAS. First, users must provide explicit,
informed consent at registration, specific to the purpose of campus
safety. Second, collected data may only be processed for declared safety
purposes and cannot be transferred to third parties or used for
advertising. Third, every user retains the right to access, correct, and
request deletion of their data, and the system must provide a mechanism
for exercising these rights. Fourth, incident records should be subject
to a defined retention policy --- a maximum of twelve months is
recommended --- after which records are deleted or irreversibly
anonymized.

For sensitive incident types such as harassment or threats, the system
provides an anonymous reporting option. The reporter's identity is
encrypted at rest and accessible only to system administrators under a
documented, access-controlled protocol; every access is recorded in the
immutable audit log. A formal data processing agreement with the
university's legal counsel is a prerequisite to production
deployment~\cite{ley1581}.

\subsection{Equity and Accessibility}

A mobile-first design must not create an equity gap for campus members
who lack a smartphone or reliable mobile data. The multi-channel
notification architecture (NFR-02) includes an SMS gateway as both a
delivery and submission fallback, ensuring that all users can send and
receive security information regardless of device access. The web portal
provides an additional access path for WiFi-connected users. All
user-facing components use Spanish as the primary language, and the
onboarding flow must be completable in under two minutes (NFR-06). The
system should not presuppose advanced digital literacy.

\subsection{Prevention of Misuse}

Three misuse vectors are explicitly addressed. Intentional false reports
are mitigated by the Verification Engine and human moderation layer
(Section~III-C). Harassment through the platform is deterred by the
immutable audit log, which records all moderation decisions and provides
accountability for every consequential action. The anonymous reporting
option includes rate limiting to prevent a single user from submitting
an excessive number of anonymous reports within a configurable time
window. These mechanisms ensure that the CSAS cannot be used to harm
the community it is designed to protect.

% =============================================================================
\section{Conclusion and Next Steps}
% =============================================================================

\subsection{Design Summary}

This document translates the empirical findings of Workshop~No.~1 into
a complete architectural blueprint for the CSAS. Each of the three core
problems identified through six independent data collection methods has
a direct and traceable design response: the reporting culture gap is
addressed by the low-friction mobile interface and community engagement
module; the non-uniform incident distribution is addressed by the
adaptive zone-weighted and time-weighted dispatch logic; and the
operational performance gaps are formalized as KPI targets in
Section~VI.

The design is guided by five systems engineering principles:
\textit{modularity} (independent microservices), \textit{scalability}
(per-component auto-scaling), \textit{maintainability} (containerized
CI/CD deployment), \textit{reliability} (redundancy and offline
queuing), and \textit{privacy} (compliance-by-design with Ley~1581
de~2012). Every major design decision is traceable to a specific
finding from Workshop~No.~1.

\subsection{Open Assumptions Requiring Validation}

\textit{Willingness to report.} The design assumes that users will
submit reports if friction is sufficiently reduced. Workshop~No.~1
survey data provides directional support (all incident-experienced
participants would use the app), but this has not been confirmed under
real operating conditions. It will be tested in the Phase~3 UAT.

\textit{Verification latency.} The design assumes that the verification
mechanism can operate within a~45-second latency budget. This depends on
implementation-specific choices not yet finalized and will be validated
through Phase~4 load testing.

\textit{Mobile connectivity coverage.} The system assumes reliable
connectivity across all campus zones. The observation phase of Workshop
No.~1 did not assess signal coverage systematically. A campus-wide
connectivity audit is recommended before Phase~1 infrastructure
planning.

\subsection{Recommendations for Workshop No.~3}

\textit{Prioritize Phase~1.} Deploying the minimum viable reporting
loop will generate real operational data to validate the sensitivity
parameters identified in Workshop~No.~1.

\textit{Conduct a connectivity audit.} Assess WiFi and mobile signal
coverage across the four zones observed during Workshop~No.~1 before
finalizing the notification delivery architecture.

\textit{Expand the student survey.} A minimum sample of~200 respondents
(stratified by faculty and time of day) is needed to validate the 35\%
incident prevalence estimate from the 20-student Workshop~No.~1 survey.

\textit{Establish a formal agreement with SIURE~UD.} A data-sharing
agreement is a prerequisite for the Section~III-D integration; it should
be initiated early, as institutional approval typically requires
significant lead time.

\textit{Engage university legal counsel.} A formal data processing
agreement under Ley~1581 de~2012 and official institutional endorsement
are both prerequisites for reaching the 60\% community adoption target
and must be pursued in parallel with Phase~1 development.

Consistent with the principle that design is iterative, this document
represents a design baseline, not a final specification. All decisions
are subject to revision as implementation proceeds and real data becomes
available.

% ── References ────────────────────────────────────────────────────────────
\begin{thebibliography}{9}
\bibitem{sierra2026}
C.~A.~Sierra,
\textit{TeamWork as Computer Engineers --- CS Collaboration Guidelines}.
Bogot\'a: Universidad Distrital Francisco Jos\'e de Caldas, 2026.

\bibitem{checkland1981}
P.~Checkland,
\textit{Systems Thinking, Systems Practice}.
New York: John Wiley \& Sons, 1981.

\bibitem{meadows2008}
D.~H.~Meadows,
\textit{Thinking in Systems: A Primer}.
White River Junction, VT: Chelsea Green Publishing, 2008.

\bibitem{iso15288}
ISO/IEC/IEEE,
\textit{Systems and Software Engineering --- System Life Cycle
Processes}, ISO/IEC/IEEE Std.~15288:2015, 2015.

\bibitem{creswell2018}
J.~W.~Creswell and J.~D.~Creswell,
\textit{Research Design: Qualitative, Quantitative, and Mixed Methods
Approaches}, 5th~ed.
Thousand Oaks, CA: SAGE Publications, 2018.

\bibitem{ley1581}
Congreso de Colombia,
\textit{Ley 1581 de 2012 --- R\'egimen General de Protecci\'on de Datos
Personales}, 2012.

\bibitem{bass2013}
L.~Bass, P.~Clements, and R.~Kazman,
\textit{Software Architecture in Practice}, 3rd~ed.
Boston: Addison-Wesley, 2013.

\bibitem{dragoni2017}
N.~Dragoni \textit{et al.},
``Microservices: Yesterday, Today, and Tomorrow,''
\textit{Present and Ulterior Software Engineering},
pp.~195--216, 2017.
\end{thebibliography}

\end{document}
